\section{Actual normalized theta points and the complete single-frame theorem}
\label{sec:normalized-frames-v050}

\status{NFR1--NFR13 construct the original \texttt{ThetaData},
\texttt{IsNormUn} and every clause of \texttt{FrameOK}, on the actual
family and the actual universal line.  They retain the original supplied
\texttt{Moduli}; HUF21 and PUN11 supply the family and its models.
The Gaussian mean is exactly one.  The three further fields of
\texttt{FramesOK}---triangular weight covariance, arithmetic-group
invariance and holomorphic readings---are not asserted here.
No principal-uniformization or whole RH/GRH obligation is closed.}

Fix the original number field $E$, CM host $H$, datum $D_w$, level $l$,
moduli record $X_o$, coefficient field $c.F$, and embedding
$\sigma:c.F\hookrightarrow\mathbb C$ satisfying \texttt{GoodEmb}.
Let
\[
 e:\{1,\ldots,g\}\simeq
       \coprod_{\rho:E\hookrightarrow\mathbb C}\{0,\ldots,D_w.\mathrm{mult}(\rho)-1\},
 \qquad \mathrm{ty}_j=\operatorname{relab}(c,\sigma)((ej)_1).
 \tag{NFR.1}
\]
Use the original $F:\texttt{FamilyUnif}$ and
$M:\texttt{ModelsOf}(F,r)$, with \texttt{ModelBase} and \texttt{NegXi}.
These are exactly the conclusions of HUF21 and PUN11 in their stated
scope, not new arithmetic hypotheses.  Fix a component $k$ and write
\[
\begin{gathered}
D=F.D_i(k),\quad \Lambda=D.\mathrm{La},\quad
\phi=M.\phi_k:\Lambda\simeq_{\mathcal O_E}M.X_k.\Lambda,\\
G=M.X_k.G_m,\quad \xi_0=M.X_k.\xi_0,\quad q=D.q.
\end{gathered}
\]
The original trace identity is
\[
 E_D(a,b)=\operatorname{Tr}_{E/\mathbb Q}
                  \bigl(-\xi_0\,\operatorname{herm}_G(\phi a,\phi b)\bigr).
 \tag{NFR.2}
\]
Every Hermitian form in this chapter is linear in its first argument:
$H(z,w)=z^tH\overline w$.  Thus the first $q$ diagonal signs of
the source's \texttt{sgn} are $+1$, and the remaining $r-q$ signs
are $-1$.  This convention is essential for (NFR.2).

\subsection{The exact frame domain and a simultaneous theta datum}

% BEGIN NFR1
\paragraph{NFR1: the original embedding and column dictionary.}
\label{claim:NFR1}
Put $\Psi=M.X_k.\Psi$, and let $\mathcal G$ be its actual frame
domain.  There is a bijection
$j\leftrightarrow(\tau,a)\in\Psi\times\{0,\ldots,r-1\}$ such that
\[
\begin{array}{c|c|c}
 &\mathrm{ty}_j& a\\ \hline
 a<q&\overline\tau&q-1-(ej)_2\\
 a\ge q&\tau&q+(ej)_2
\end{array}                                                     \tag{NFR.3}
\]
For a frame $x=(X_\tau)_{\tau\in\Psi}$ and $v\in E^r$, its
literal \texttt{perFr} coordinates are
\[
 p_x(v)_j=
 \begin{cases}
  (\operatorname{col}_aX_\tau)^tG_\tau\overline{\tau(v)},&a<q,\\
  (X_\tau^{-1}\tau(v))_a,&a\ge q.
 \end{cases}                                                    \tag{NFR.4}
\]
In particular $P_x(a)=p_x(\phi a)$ defines the unique additive
map with the source's \texttt{IsPerFr} formula.

\begin{proof}
The map $\operatorname{relab}(c,\sigma)$ permutes the embeddings
of $E$: it comes from applying the coefficient-field embedding to
the chosen realization of every embedding in its splitting field.
Its inverse is the inverse field automorphism on that splitting field.
It commutes with conjugation, since conjugation on a CM field is
the unique involution over its maximal real subfield.  The
\texttt{GoodEmb} equation says it preserves the test
$\rho(H.\mathrm{rt})=H.\mathrm{rt}_0$.
Consequently the type has multiplicity $D_w.p$ at each member of
$\Psi$ and $D_w.q$ at its conjugate.  Since
$r=D_w.p+D_w.q$ and $q=D_w.q$, send the former letters to
columns $q,q+1,\ldots,r-1$ and the latter, in reverse order,
to $q-1,q-2,\ldots,0$.  This is (NFR.3), including the empty
blocks when $p$ or $q$ is zero.

Substitution in the definition of \texttt{perFr} gives (NFR.4).
The two cases are additive in $v$; multiplication of $v$ by
$a\in\mathcal O_E$ multiplies the first case by $\overline\tau(a)$
and the second by $\tau(a)$.  These are precisely $\mathrm{ty}_j(a)$.
The inclusion of $M.X_k.\Lambda$ in $E^r$ composed with $\phi$
is injective, has that exact range, and spans $E^r$.
The trace, conjugation and signature data are the corresponding
\texttt{ModelBase}, \texttt{Data} and \texttt{NegXi} fields.
No different rational lattice or reordered type is substituted.
\end{proof}
% END NFR1

% BEGIN NFR2
\paragraph{NFR2: these are holomorphic periods on an open connected domain.}
\label{claim:NFR2}
For every $x\in\mathcal G$, the map $P_x$ of NFR1 satisfies the
entire original \texttt{IsPeriod} predicate.  Each function
$x\mapsto P_x(a)$ is holomorphic on $\mathcal G$, and
$\mathcal G$ is open and connected.

\begin{proof}
Write $A_\tau=X_\tau^tG_\tau\overline{X_\tau}$.
The defining inequalities of $\mathcal G$ say that its leading
minors are nonzero with successive signs
$s_a=1$ for $a<q$, $s_a=-1$ otherwise.
Hermitian elimination gives an upper unitriangular $U_\tau$ with
\[
 U_\tau^tA_\tau\overline{U_\tau}
       =\operatorname{diag}(d_{\tau,a}),\qquad s_a d_{\tau,a}>0.
 \tag{NFR.5}
\]
Here one successively replaces column $a$ by itself minus its
projections onto earlier orthogonal columns, and makes the
conjugate row operations.  The $a$th diagonal entry is the ratio
of consecutive leading determinants.  This proves the stated signs
and also proves that the algorithm never divides by zero.
Rescaling each column by $|d_{\tau,a}|^{-1/2}$ gives $X_{0,\tau}$
with Gram matrix $\operatorname{diag}(s_a)$ and
$X_\tau=X_{0,\tau}b_\tau$ for an invertible upper triangular $b_\tau$.

In an orthonormal frame, put $c=X_{0,\tau}^{-1}\tau(v)$.
Formula (NFR.4) is $\overline c_a$ for $a<q$ and $c_a$ otherwise.
It is therefore a real linear isomorphism.  For arbitrary $b_\tau$,
the period vector changes, in the unreversed column order, by
\[
 Q_\tau=
 \begin{pmatrix} (b_{\tau,++})^t&0\\0&(b_{\tau,--})^{-1}\end{pmatrix}.
 \tag{NFR.6}
\]
Indeed an early column involves only earlier columns; a late row
of $b_\tau^{-1}$ involves only later rows.  Thus (NFR.6) is
invertible and the whole period map remains a real linear isomorphism.

For $\xi=-\xi_0$, write $\tau(\xi)=i\alpha_\tau$,
$\alpha_\tau=|\tau(\xi)|>0$.
The trace identity, pairing conjugate embeddings, becomes
\[
 E_D(a,b)=-2\sum_{\tau\in\Psi}\alpha_\tau
          \operatorname{Im}\operatorname{herm}_{G_\tau}
                         (\tau\phi a,\tau\phi b).
 \tag{NFR.7}
\]
In the orthonormal coordinates this is the imaginary part of the
positive diagonal form with entries $2\alpha_\tau$ on the
coordinates (NFR.4).  Conjugating that form by $Q^{-1}$ supplies
the positive Hermitian form for arbitrary $x$.
The image of the actual $\mathcal O_E$-lattice is discrete and
spans the full real space: after choosing an $E$-basis, it lies
between two nonzero integral multiples of $\mathcal O_E^r$,
whose Minkowski embedding is a full real lattice.  The real
isomorphism just proved preserves these properties.
Together with NFR1 this checks every \texttt{IsPeriod} clause.

The negative-type coordinates in (NFR.4) are complex linear in
the entries of $X_\tau$; the other coordinates are rational with
denominator $\det X_\tau$.  They are holomorphic on the domain.
Its definition is a finite conjunction of open inequalities in
continuous minors and nonzero determinants, so it is open.

For connectedness, (NFR.5) expresses every frame as an
orthonormal frame times an upper triangular matrix with positive
real diagonal.  The latter set is connected.  The orthonormal
frames for a fixed $G_\tau$ form a translate of $U(q,r-q)$.
This group is connected: a positive $q$-plane for the standard
form is uniquely the graph of a matrix $Z$ with $Z^*Z<I$;
the matrix ball is convex, and the positive square roots of
$I-Z^*Z$ and $I-ZZ^*$ give a continuous orthonormal graph frame.
The remaining choices are $U(q)\times U(r-q)$, each connected
by diagonalizing a unitary matrix and varying its eigenvalue
arguments to zero.  This proves connectedness of each factor,
then of their finite product and of $\mathcal G$.
The argument also includes $q=0$, $q=r$, and $r=0$.
\end{proof}
% END NFR2

% BEGIN NFR3
\paragraph{NFR3: an integral symplectic basis with the exact modulus eight.}
\label{claim:NFR3}
There is a basis
$\mathrm{bas}:\mathbb Z^g\oplus\mathbb Z^g\simeq\Lambda$ such that
\[
 E_D(\mathrm{bas}(u,v),\mathrm{bas}(u',v'))
                      =-u\cdot v'+v\cdot u',
 \tag{NFR.8}
\]
and, for the actual host element $w_0$,
\[
 w_0^4\mathrm{bas}(e_i,0)\equiv\mathrm{bas}(e_i,0)\pmod{8\Lambda},
 \qquad w_0^4\mathrm{bas}(0,e_i)\equiv0\pmod{8\Lambda}.
 \tag{NFR.9}
\]

\begin{proof}
The finite projective $\mathcal O_E$-module $\Lambda$ is finite
and torsion free over $\mathbb Z$, hence free.  Its rank is
$[E:\mathbb Q]r=2g$ by NFR1.  The form $E_D$ is unimodular
and alternating.  An integral symplectic basis exists by repeatedly
choosing a primitive vector, using unimodularity to pair it to one,
and splitting off the resulting hyperbolic plane.

The host equations $2w_0=1+\mathrm{rt}$,
$\mathrm{rt}^2=-d$ and $d\equiv7\pmod {16}$ imply, with
$m=2\lfloor d/16\rfloor+1$,
\[
 w_0^2=w_0-2m,\quad \overline{w_0}=1-w_0,\quad
 w_0^4+(1-w_0)^4\equiv1,\quad
 w_0^4(1-w_0)^4\equiv0\pmod8.
 \tag{NFR.10}
\]
Thus multiplication by $w_0^4$ on $\Lambda/8\Lambda$ is an
idempotent $e_+$; its adjoint for $E_D$ is $1-e_+$.
Its image and kernel are complementary isotropic direct summands.
The perfect form pairs them perfectly, so both are free of rank
$g$ over the local ring $R=\mathbb Z/8\mathbb Z$.
Choose a basis of the image and its dual basis in the kernel,
with the sign in (NFR.8).  This gives a symplectic basis modulo eight.

We justify lifting this basis integrally.  The reduction map
$\operatorname{Sp}_{2g}(\mathbb Z)\to\operatorname{Sp}_{2g}(R)$
is surjective by the following elementary elimination.
The symmetric shears
$\left(\begin{smallmatrix}I&S\\0&I\end{smallmatrix}\right)$,
$\left(\begin{smallmatrix}I&0\\S&I\end{smallmatrix}\right)$,
the simultaneous elementary changes
$\operatorname{diag}(C,C^{-t})$, and signed interchanges of
hyperbolic pairs all lift when $C$ is elementary.
The first column of a symplectic matrix is unimodular; over the
local ring $R$ it has a unit entry.  Move that entry to its first
upper coordinate.  A lower symmetric shear kills all lower
coordinates: choose its off-diagonal entries in the first column
to kill coordinates $2,\ldots,g$, and then its $(1,1)$ entry
to kill the first.  Elementary changes kill the remaining upper
coordinates.  A diagonal scaling on the first hyperbolic pair
normalizes the surviving unit.  This scaling also lifts, because
for every unit $u$ the product of the three two-dimensional shears
$U(u)L(-u^{-1})U(u)$, followed by the signed interchange, is
$\operatorname{diag}(u,u^{-1})$; lift the entries of the shears.

The first column is now $e_1$.  Its symplectic partner has lower
first coordinate one.  An upper symmetric shear fixing $e_1$
kills its upper coordinates.  An elementary change fixing $e_1$
then kills its remaining lower coordinates.  The other columns,
orthogonal to this pair, have zero first-pair coordinates; induction
finishes the reduction.  Reversing these lifted operations lifts
the given symplectic matrix.  Apply such a lift to the initial
integral symplectic basis.  Its reduction is the chosen image/kernel
basis, so (NFR.9) holds in the literal lattice, proving the claim.
\end{proof}
% END NFR3

% BEGIN NFR4
\paragraph{NFR4: the whole original \texttt{ThetaData} exists.}
\label{claim:NFR4}
For every component $k$ there is an actual
$T:\texttt{ThetaData}(e,F,M,k)$, with its original total period
function, coordinate isomorphism, adapted basis, normalizing
matrices, fixed half-characteristic and universal-line frames.

\begin{proof}
Take $T.Pf(x)=P_x$ from NFR1--NFR2 and the basis of NFR3.
Choose a complex linear coordinate isomorphism
$\mathrm{co}:\mathbb C^m\simeq\prod_{\tau\in\Psi}M_r(\mathbb C)$;
the preimage $B=\mathrm{co}^{-1}(\mathcal G)$ is open and connected.
Let $A(b),B_1(b)$ be the column matrices of the first and second
halves of $P_{\mathrm{co}(b)}\mathrm{bas}$.
The first real half-lattice spans a real isotropic $g$-plane $W$.
Its complex span is all of $\mathbb C^g$: otherwise $W\cap iW$
would contain $iw\ne0$ with $w\in W$, and
$\operatorname{Im}H(iw,w)=H(w,w)>0$ would contradict isotropy.
Thus $A$ is invertible.  Set $h=A^{-1}$ and $\tau=A^{-1}B_1$.

For clarity, the Riemann relations follow directly here.
The form $H'=A^tH\overline A$ is real symmetric and positive,
since its imaginary part vanishes on the first half-basis.
The mixed pairing in (NFR.8) says
$H'\operatorname{Im}\tau=I$.  Hence $Y=\operatorname{Im}\tau$
is positive and $H'=Y^{-1}$.  The imaginary part on the second
half-basis then says $\operatorname{Re}\tau$ is symmetric.
This gives every normalization equation in $T.\mathrm{norm}$.
The matrices $h,\tau$ are holomorphic because the period entries
are holomorphic and $\det A\ne0$.  Define them arbitrarily outside
$B$; all their asserted properties have the original domain guard.

Apply the actual $F.\mathrm{theta}$ field, supplied in full by
HUF20--HUF21, on this open connected base, with this prescribed
basis, sign $-1$, and these matrices.  It gives two constant
half-characteristics $c_1,c_2\in\{0,1/2\}^g$ and the actual
functions $fr_j(b,z)$ in the charts of the universal rigidified
line.  Its five clauses give exactly $T.fr_{\rm hol}$,
$T.fr_{\rm ne}$, $T.fr_{\rm coc}$, $T.fr_{\rm rig}$ and
$T.fr_{\rm sec}$.  Together with NFR1--NFR3 these fill every
field of \texttt{ThetaData}, including its modulo-eight
\texttt{adapt} field.  This proves the full source
\texttt{nonempty\_thetaData} conclusion.
\end{proof}
% END NFR4

\subsection{The actual algebraic section and its normalization}

% BEGIN NFR5
\paragraph{NFR5: theta convergence, nonvanishing and the section line.}
\label{claim:NFR5}
Let $\tau=X+iY$ be symmetric with $Y>0$, and let $a,b$ be real
half-characteristics.  With the source convention,
\[
 \theta[a,b](\tau,w)=\sum_{n\in\mathbb Z^g}
   \exp\{\pi i(n+a)^t\tau(n+a)+2\pi i(n+a)^t(w+b)\}.
 \tag{NFR.11}
\]
It is entire in $w$, is not identically zero, and
\[
 \theta(w+m+\tau n)=J_{m,n}(w)\theta(w),\quad
 J_{m,n}(w)=e^{\pi i(2a\cdot m-2n\cdot b-n^t\tau n-2n\cdot w)}.
 \tag{NFR.12}
\]
For a fiber of the $T$ in NFR4, its actual algebraic section
$s_0$ from $T.fr_{\rm sec}$ is nonzero and spans the whole space
of algebraic sections of the actual principal line.

\begin{proof}
On a compact set of $w$, the norm of a summand is at most
$\exp(-c|n|^2+C|n|+C)$ for $c>0$ depending on the least
eigenvalue of $Y$.  The same bound with a polynomial factor holds
after any derivative.  Thus the series and its derivatives converge
normally.  Translation by $m$ gives its characteristic factor;
translation by $\tau n$ follows by shifting the summation index.
This proves (NFR.12) with the stated sign.
Multiplying by $e^{-2\pi ia\cdot w}$ and integrating against
$e^{-2\pi ik\cdot u}$ for $u\in[0,1]^g$ extracts the nonzero
coefficient
$e^{\pi i(k+a)^t\tau(k+a)+2\pi i(k+a)\cdot b}$.
Hence the function is not zero.

We spell out why this yields the actual one-dimensional section
space, without adding a section-rank assumption.
Fix the fiber and pull an arbitrary algebraic section back by its
exponential.  Division by the nowhere-zero $T.fr_j$ gives a
single entire function $f$ on $\mathbb C^g$: the chart transition
identities ensure agreement on overlaps.  The section $s_0$
has scalar $\theta(hz)$.
Since exponential points $z$ and $z+\lambda$ are equal, its
values and (NFR.12) imply
$fr_j(z)=J_{m,n}(hz)fr_j(z+\lambda)$ whenever
$h\lambda=m+\tau n$.
Initially cancel $\theta(hz)$ on its nonzero set.  This set is
dense: if an entire function vanished on a ball, its restriction
to every complex line through that ball would vanish identically.
Continuity extends the identity across all its zeros.
The same identity applied to $f$ shows that $F(w)=f(h^{-1}w)$
satisfies (NFR.12).

Every entire function with the $\mathbb Z^g$-translation law
has the normally convergent Fourier expansion
$F(w)=\sum_{k\in\mathbb Z^g}c_k e^{2\pi i(k+a)\cdot w}$.
One may obtain it as the Laurent expansion on $(\mathbb C^*)^g$
after multiplying by $e^{-2\pi ia\cdot w}$; equivalently its
coefficients are the preceding cube integrals, independent of
the imaginary slice by Cauchy's theorem.
Substitution of the $\tau n$-translation law yields
\[
 c_{k+n}=c_k\exp\{\pi i((k+n+a)^t\tau(k+n+a)
                  -(k+a)^t\tau(k+a))+2\pi in\cdot b\}.
 \tag{NFR.13}
\]
All coefficients are therefore one scalar multiple of the
coefficients in (NFR.11), and $F$ is that multiple of $\theta$.
It follows that every algebraic section is a scalar multiple of
$s_0$: equality on complex points is equality of regular functions
in each chart, because the fiber is smooth, reduced and of finite
type over $\mathbb C$.  Nonvanishing of $\theta$ and of a chart
frame proves $s_0\ne0$, and hence uniqueness of the scalar.

This proves exactly \texttt{SecGen} for every nonzero multiple
of $s_0$.  It agrees with TDR2's section-line description wherever
\texttt{SecBaseChangeSt} is supplied.  The argument here only uses
the already supplied actual $\texttt{ThetaData}$ and the complex
fiber, so it also proves the original conclusions stated with
\texttt{SecTwistSt} as a premise, without asserting that premise
for any other base or line bundle.
\end{proof}
% END NFR5

% BEGIN NFR6
\paragraph{NFR6: the whole actual normalized-point constructor.}
\label{claim:NFR6}
For every $b\in T.co^{-1}(\mathcal G)$ and every
$\kappa\in\mathbb C^\times$, there is an actual
$y:\texttt{FrPt}(X_o,c,CAt(\sigma))$ satisfying
$\texttt{IsThetaPt}(T,b,\kappa,y)$.
Consequently there is a single function
$un:\mathcal G\to\texttt{FrPt}(X_o,c,CAt(\sigma))$
satisfying the full original \texttt{IsNormUn} predicate:
for every $x\in\mathcal G$ some $\kappa_x$ satisfies
\[
 \kappa_x^2=C_0\det h(T.co^{-1}x),\qquad
 C_0=\sqrt{\frac{2^g}{P_0}},\qquad
 P_0=\prod_{j=1}^g2|\mathrm{ty}_j(\xi_0)|,
 \tag{NFR.14}
\]
and $\texttt{IsThetaPt}(T,T.co^{-1}x,\kappa_x,un(x))$.

\begin{proof}
Put $y.x=F.u_k(P_{T.co(b)})$.  Its coefficient map is the one in
$F.\mathrm{over}_\sigma$.  Differentiate the actual exponential
in the coordinate directions $e_j$.  The two $F.\mathrm{ex}_{\rm diff}$
clauses make this differential a bijection from $\mathbb C^g$
onto the tangent space at the identity.  Additivity of the
exponential makes it complex linear; $F.\mathrm{ex}_{\rm lin}$
makes direction $j$ an eigenvector with eigenvalue
$\mathrm{ty}_j(a)$.  Under the original embedding (NFR.1) these
are precisely the ordered eigenbases in \texttt{FrPt}.
HUF1 identifies these analytic derivatives with the literal
dual-number predicate \texttt{IsDerivAt0}, and TDR1 identifies
linear combinations with its literal \texttt{IsCombAt}.

Transport the actual $s_0$ of NFR5 along the canonical isomorphism
between the two pullback models of the same fiber, and set
$y.s=\kappa s_0$.  Pullback does not change a chart value at a
point with the two specified projections; multiplication by
$\kappa$ multiplies that value by $\kappa$.
Thus its value is exactly
$\kappa\theta(T.ta(b),T.c_1,T.c_2,T.h(b)z)T.fr_j(b,z)$.
NFR5 proves the literal \texttt{SecGen} condition.  These data
fill every field of $y$ and every \texttt{IsThetaPt} clause.

Each factor in $P_0$ is positive, since $\xi_0\ne0$ and field
embeddings are injective.  Thus $C_0>0$.  The determinant in
(NFR.14) is nonzero.  Choose a square root of its nonzero
product with $C_0$, and apply the preceding constructor.
Choice at every $x$ produces one function $un$.
No continuity of this square-root choice is required or claimed;
\texttt{IsNormUn} has precisely this pointwise existential
quantifier.  This proves the entire original
\texttt{exists\_isNormUn} conclusion, and not a prescribed-sign
or simply-connected-base variant.
\end{proof}
% END NFR6

\subsection{The two determinant identities and the Gaussian mean}

% BEGIN NFR7
\paragraph{NFR7: the normalized period matrix determines the Hermitian form.}
\label{claim:NFR7}
For any original $T$, put $P=T.Pf(T.co(b))$, $h=T.h(b)$,
$\tau=T.ta(b)$ and $Y=\operatorname{Im}\tau$.
If $H$ is a matrix with
$\operatorname{Im}H(Pa,Pb)=E_D(a,b)$, then
\[
 H=h^tY^{-1}\overline h,
 \qquad \det H\det Y=|\det h|^2.
 \tag{NFR.15}
\]

\begin{proof}
Let $A,B$ be the two column matrices of the prescribed integral
basis.  Then $hA=I$, $hB=\tau$, so $B=A\tau$.
Put $R=A^tH\overline A$.  The first-half pairing says
$\operatorname{Im}R=0$, while the mixed pairing says
$\operatorname{Im}(R\overline\tau)=-I$.
Thus $RY=I$ and $R=Y^{-1}$.  Conjugating back gives the first
identity.  Taking determinants and using
$\det\overline h=\overline{\det h}$ gives the second.
This argument does not silently replace $H$ by a different
polarization form; it determines the given form itself.
\end{proof}
% END NFR7

% BEGIN NFR8
\paragraph{NFR8: the exact determinant for unit-minor frames.}
\label{claim:NFR8}
Suppose $x\in\mathcal G$ has the source normalization
$|\operatorname{dmin}(G,x,\tau,i)|=1$ for every $\tau,i$.
Then every $H$ with the original imaginary period pairing satisfies
\[
                     \det H=P_0.                     \tag{NFR.16}
\]
This is the full original \texttt{UnitDet.det\_eq\_prod} conclusion.

\begin{proof}
Every leading Gram determinant is real, has its prescribed sign,
and has absolute value one.  Hence it equals the product of its
first $i$ signs.  In the unitriangular elimination (NFR.5), the
successive determinant ratios are therefore exactly $s_a$.
No diagonal rescaling is needed: $X=X_0b$ with each $b_\tau$
upper unitriangular.  Formula (NFR.6) shows that
$P_x=Q P_{x_0}$ with $\det Q=1$; in particular $|\det Q|=1$.
All identities hold after the exact permutation (NFR.3).

At $x_0$, (NFR.7) says the imaginary pairing is given by
\[
              H_0=\operatorname{diag}
                           (2|\mathrm{ty}_j(\xi_0)|)_{j=1}^g.
\]
The form $Q^tH\overline Q$ has that same imaginary pairing on
the period lattice of $x_0$.  Two complex sesquilinear forms
whose imaginary parts agree on a real spanning set agree
everywhere: extend by real bilinearity to all vectors, then
apply the equality also to $iz,w$ to recover their real parts.
Thus $Q^tH\overline Q=H_0$.
Taking determinants gives $|\det Q|^2\det H=\det H_0=P_0$,
as asserted.  The reverse order of the first $q$ columns and
the sign $-\xi_0$ in the trace were used in (NFR.7); dropping
either would invalidate the positive diagonal formula.
\end{proof}
% END NFR8

% BEGIN NFR9
\paragraph{NFR9: an actual fundamental domain and its exact theta mean.}
\label{claim:NFR9}
For $h\in\operatorname{GL}_g(\mathbb C)$, $\tau=X+iY$ with
$Y>0$ symmetric, and the exact lattice
$\mathcal L=h^{-1}(\mathbb Z^g+\tau\mathbb Z^g)$, define
\[
 F_d=\{h^{-1}(u+\tau v):u,v\in[0,1)^g\}.
\]
It is an additive fundamental domain for $\mathcal L$, has finite
positive volume, and for real $a,b$ (in particular the source's
half-characteristics),
\[
 \frac1{\operatorname{vol}F_d}\int_{F_d}
   |\theta[a,b](\tau,hz)|^2
    e^{-2\pi\operatorname{Im}(hz)^tY^{-1}\operatorname{Im}(hz)},dz
                       =\frac1{\sqrt{2^g\det Y}}.
 \tag{NFR.17}
\]

\begin{proof}
The real linear map $(u,v)\mapsto h^{-1}(u+\tau v)$ is invertible:
the imaginary part recovers $v$ using $Y^{-1}$, and then the
real part recovers $u$.  Its image of the half-open unit cube
tiles the whole space by the exact image of $\mathbb Z^{2g}$.
Its constant nonzero Jacobian is both $\operatorname{vol}F_d$
and the change-of-variables factor, so it cancels in the mean.

At $w=u+\tau v$, integrate first in $u$.
The normally convergent series from NFR5 permits termwise
integration of its square on the compact cube.
Fourier orthogonality eliminates every pair of different indices.
After multiplying by $e^{-2\pi v^tYv}$, the remaining summand is
\[
                       e^{-2\pi(n+a+v)^tY(n+a+v)}.
\]
All these terms are nonnegative; Tonelli's theorem permits their
integration and summation in either order.  The cubes
$n+a+[0,1)^g$ partition $\mathbb R^g$ up to their null boundaries.
Thus the double cube integral is
$\int_{\mathbb R^g}e^{-2\pi t^tYt}\,dt$.
Diagonalize the positive real matrix $Y$ orthogonally and scale
each coordinate; the product of the one-dimensional Gaussian
integrals is $2^{-g/2}(\det Y)^{-1/2}$.
For completeness, $\int_{\mathbb R}e^{-\pi t^2}\,dt=1$
follows by squaring the integral and using polar coordinates;
scaling gives the factors just used.
This proves (NFR.17), with the original Lebesgue measure and
normalization by the volume of this very domain.
\end{proof}
% END NFR9

% BEGIN NFR10
\paragraph{NFR10: the holomorphic gauge and its unit lattice multiplier.}
\label{claim:NFR10}
With $H=h^tY^{-1}\overline h$, set
\[
 q_0(z)=\frac\pi2(hz)^tY^{-1}(hz),\qquad
 \Theta(z)=e^{q_0(z)}\theta[a,b](\tau,hz).
\]
For $h\lambda=m+\tau n$ define
\[
 \chi(\lambda)=e^{\pi i(m\cdot n+2a\cdot m-2n\cdot b)}.
\]
Then $|\chi(\lambda)|=1$ and
\[
 \Theta(z+\lambda)=\chi(\lambda)
       e^{\pi H(z,\lambda)+\pi H(\lambda,\lambda)/2}\Theta(z),
 \tag{NFR.18}
\]
\[
 |e^{q_0(z)}|^2e^{-\pi H(z,z)}
             =e^{-2\pi\operatorname{Im}(hz)^tY^{-1}\operatorname{Im}(hz)}.
 \tag{NFR.19}
\]
Here $H(z,z)$ is real.

\begin{proof}
Write $w=hz$, $\mu=m+\tau n$ and $S=Y^{-1}$.
The relation $\overline\mu=\mu-2iYn$ gives
$w^tS\overline\mu=w^tS\mu-2iw\cdot n$.
Expand $q_0(z+\lambda)-q_0(z)$ and insert (NFR.12).
The $w$-dependent terms become $\pi w^tS\overline\mu$.
The constant difference from
$\pi\mu^tS\overline\mu/2$ is
$\pi i(m\cdot n+2a\cdot m-2n\cdot b)$.
This proves (NFR.18); its exponent is purely imaginary, proving
the unit norm.  For (NFR.19), write $w=u+iv$.
Then $\operatorname{Re}(w^tSw)=u^tSu-v^tSv$ and
$w^tS\overline w=u^tSu+v^tSv$.
Subtracting gives $-2v^tSv$, exactly the exponent in (NFR.19).
\end{proof}
% END NFR10

% BEGIN NFR11
\paragraph{NFR11: normalized square roots give mean exactly one.}
\label{claim:NFR11}
At a unit-minor frame $x$, let $\kappa$ be any root in (NFR.14).
Then
\[
 \frac{|\kappa|^2}{\sqrt{2^g\det Y}}=1,\qquad
 \frac1{\operatorname{vol}F_d}\int_{F_d}
  |\kappa\Theta(z)|^2e^{-\pi H(z,z)},dz=1.            \tag{NFR.20}
\]

\begin{proof}
NFR7--NFR8 give $P_0\det Y=|\det h|^2$.
Taking norms of $\kappa^2=C_0\det h$ gives
$|\kappa|^2=C_0|\det h|$.
Every factor is positive, so
\[
 C_0|\det h|=\sqrt{(2^g/P_0)|\det h|^2}
                        =\sqrt{2^g\det Y}.
\]
This proves the first equality, independently of the chosen sign
of $\kappa$.  Multiply (NFR.17) by $|\kappa|^2$ and use
(NFR.19) to obtain the second.
\end{proof}
% END NFR11

% BEGIN NFR12
\paragraph{NFR12: every literal clause of the original \texttt{FrameOK}.}
\label{claim:NFR12}
For every original $T:\texttt{ThetaData}(e,F,M,k)$ and every
function $un$ satisfying \texttt{IsNormUn}, one has
\[
                  \forall x\in\mathcal G,\quad
                     \texttt{FrameOK}(e,F,M,k,x,un(x)).
 \tag{NFR.21}
\]
The statement retains exactly \texttt{GoodEmb}, \texttt{ModelBase}
and \texttt{NegXi}.  In particular it proves the original
\texttt{frameOK\_of\_isNormUn} theorem with its printed
\texttt{SecTwistSt} antecedent, which the proof of this implication
does not need.

\begin{proof}
Fix $x$ and the $\kappa$ supplied by \texttt{IsNormUn}; put
$b=T.co^{-1}x$ and $P=T.Pf(x)$.
The period formula, \texttt{IsPeriod}, equality of the moduli
point, and all prescribed derivatives are respectively $T.per$,
$T.isPer$ and the first two \texttt{IsThetaPt} clauses.
This gives the unconditional part of \texttt{FrameOK}.

Now impose precisely its unit-minor antecedent.
Choose the positive Hermitian $H$ already supplied by
\texttt{IsPeriod}.  NFR7 identifies it, and NFR8 computes its
determinant.  The basis normalization in $T$ identifies its exact
period subgroup with $h^{-1}(\mathbb Z^g+\tau\mathbb Z^g)$:
write any lattice vector in the integral basis for one direction;
for the converse use that same basis and invertibility of $h$.
Thus the $F_d$ in NFR9 is a fundamental domain for $P.\mathrm{range}$,
as required by the source, not a finite-index sublattice.
Use the source chart frames and scalar
\[
 \widetilde{fr}_j(z)=T.fr_j(b,z)e^{-q_0(z)},\qquad
 \widetilde\theta(z)=\kappa e^{q_0(z)}\theta(T.ta(b),T.c_1,T.c_2,T.h(b)z).
 \tag{NFR.22}
\]
Define $\chi$ on the period lattice by NFR10 and arbitrarily
off it.  Uniqueness of $m,n$ follows from $Y>0$.

The matrix $H$ is Hermitian and positive and has exactly the
original imaginary period pairing.  Each $\widetilde{fr}_j$ is
holomorphic and nonzero on the exact inverse image of its
algebraic chart.  Multiplication by the common exponential
preserves every transition equation.  Since $q_0(0)=0$, its
value at zero is exactly the given rigidification value.
These verify the first seven displayed frame clauses.

For the translation clause, let $\lambda=Pa$.
The exponential satisfies $\mathrm{ex}(z+\lambda)=\mathrm{ex}(z)$,
using its actual kernel and group law.  Thus the scalar values
of the algebraic section $s_0$ at these two presentations of
the same fiber-product point agree.  Combining that equality
with (NFR.18) yields
\[
 \widetilde{fr}_j(z)=\chi(\lambda)
    e^{\pi H(z,\lambda)+\pi H(\lambda,\lambda)/2}
                           \widetilde{fr}_j(z+\lambda).
 \tag{NFR.23}
\]
Cancellation first proves this away from the zeros of $\Theta$.
Those zeros have empty interior by NFR5, so continuity extends
it on the whole open chart.  The translated chart is the same
open set because the exponential points are equal.  Thus this
extension is legitimate even where the theta section vanishes.
NFR10 also gives $|\chi(Pa)|=1$.

The last \texttt{IsThetaPt} clause and
$e^{q_0(z)}e^{-q_0(z)}=1$ give, for every original chart and
every fiber-product point with the two prescribed projections,
the exact identity
$\operatorname{val}(un(x).s)=\widetilde\theta(z)
\widetilde{fr}_j(z)$.  Finally NFR9 gives the literal additive
fundamental-domain condition and NFR11 gives its mean equal to
one.  These are all remaining clauses of \texttt{FrameOK}.
Nothing in this argument changes the section of $un(x)$ or its
ordered tangent frame.
\end{proof}
% END NFR12

% BEGIN NFR13
\paragraph{NFR13: simultaneous actual choices on the finite family.}
\label{claim:NFR13}
Under the original input scope above, choose one $T_k$ for each
component and one function $un_k$ such that
\[
  \forall k,\quad\texttt{IsNormUn}(T_k,un_k),\qquad
  \forall k\ \forall x\in M.X_k.\mathrm{dom},\quad
                       \texttt{FrameOK}(e,F,M,k,x,un_k(x)).
 \tag{NFR.24}
\]
With the existing \texttt{SecBaseChangeSt} premise of TDR13,
every one of these very framed points is \texttt{FromTest}.

\begin{proof}
NFR4 supplies each $T_k$, then NFR6 supplies its $un_k$.
Choose them once for the finite component set, before the
universal frame variable.  NFR12 applies to every chosen
$un_k$; it is valid for any choices of the pointwise square roots.
TDR13 is universal over actual complex \texttt{FrPt} values,
so apply it directly to $un_k(x)$, preserving its moduli point,
its ordered eigenbasis and its actual generating section.
The finitely presented reduced test ring and its frame may depend
on $k,x$, exactly as the original \texttt{FromTest} permits.
\end{proof}

The full original \texttt{FramesOK} is a conjunction of four
fields.  Equation (NFR.24) supplies its first field, on one
simultaneous family.  Its other three fields need the
transformation of theta squares under the exact integral
level group, the resulting triangular character identity,
and holomorphic evaluation of the original algebraic section
rules.  They are not consequences of choosing a square root
or of the mean-one calculation.  Likewise the independent
\texttt{CoverTri}, \texttt{FibTri} and complete \texttt{IsUnf}
assembly remain open in this packet.
% END NFR13

\paragraph{Source map and inherited scope.}
The exact original contracts and source ranges are recorded in
this component's \texttt{contracts.json}, \texttt{provenance.json}
and \texttt{source\_manifest.json}.
The main source bodies are:\par
\noindent\texttt{NormConst.lean};\par
\noindent\texttt{NonemptyThetaDataFlagDomainModels.lean};\par
\noindent\texttt{ExistsThetaPt.lean};\par
\noindent\texttt{ExistsNormalisedFramedPointOnEachComponent.lean};\par
\noindent\texttt{HermEqOfNormalised.lean};
\texttt{EqOfImHermFormEq.lean};\par
\noindent\texttt{ExistsFundamentalDomainMeanNormSqThetaChar.lean};\par
\noindent\texttt{ExpMulThetaCharAutomorphy.lean};\par
\noindent\texttt{NormalisedPointsSatisfySingleFrame.lean}.\par
The proofs use HUF1, HUF20--HUF21, PUN11 and the exact TDR
frame dictionary as indicated; their source hashes are separately
bound.  No assertion about existence of the original representing
moduli scheme or the remaining global frame laws is added.
